\BeleskeID{09_1-s004}
% element: 09_1-s004:n01
Pri prelasku iz triodne oblasti u zasićenje očekuje se isti rezultat za struju. U nemodifikovanim izrazima faktor modulacije dužine kanala pojavljuje se samo u zasićenju, što može dati diskontinuitet. Zamenom graničnog uslova dobija se
\[I_{Dn}=\frac{k_n}{2}(V_{GSn}-V_{Tn})^2(1+\lambda_nV_{DSn})=\frac{k_n}{2}V_{DSn}^2(1+\lambda_nV_{DSn}),\]
\[I_{Dn}=\frac{k_n}{2}\bigl(2V_{DSn}(V_{GSn}-V_{Tn})-V_{DSn}^2\bigr)=\frac{k_n}{2}V_{DSn}^2.\]
Ovo je posledica nesavršenosti modela, slično grubim modelima bipolarnih tranzistora. Jedan pristup menja faktor u izrazu za zasićenje, a drugi isti faktor modulacije uključuje i u triodnoj oblasti. Pri daljoj analizi koristićemo nemodifikovane izraze; u velikom broju slučajeva uticaj modulacije zanemarujemo.

